\documentclass[11pt,letterpaper]{article}
\usepackage[margin=1in]{geometry}
\usepackage{amsmath,amssymb,amsthm}
\usepackage[T1]{fontenc}
\usepackage{lmodern}
\usepackage[hidelinks]{hyperref}
\usepackage{microtype}
\newtheorem{theorem}{Theorem}
\newtheorem{lemma}{Lemma}
\newcommand{\R}{\mathbb{R}}
\newcommand{\N}{\mathbb{N}}
\newcommand{\vol}{\operatorname{vol}}
\newcommand{\conv}{\operatorname{conv}}
\newcommand{\code}[1]{\texttt{\small #1}}
\title{A discreteness obstruction to the printed polytope approximation rate}
\author{C0ldSmi1e\\\small Conjecture 00000001451}
\date{6 October 2026}
\begin{document}
\maketitle
\begin{abstract}
The conjecture assigns the optimal facet count the rate
$c_d\varepsilon^{(d-1)/2}$. Both source languages print this positive
exponent. In dimension two, no natural number can lie between two fixed
positive multiples of $\sqrt\varepsilon$ for every sufficiently small
positive $\varepsilon$. Consequently the conjecture fails as an exact
formula, a positive leading asymptotic, and even a positive two-sided
order estimate. The Lean formalization retains the actual Euclidean
volume approximation condition, geometric facets, and attained minimum.
\end{abstract}

\section{Statement and scope}
Let $B_d$ be the closed unit ball in Euclidean $\R^d$. The source calls a
polytope $P$ an $\varepsilon$-approximation when
\[
 \vol(B_d\mathbin{\triangle}P)\leq\varepsilon\vol(B_d).
\]
It states that the number of facets of the minimal approximating
polytope is $c_d\varepsilon^{(d-1)/2}$, with an explicit exact constant
(the constant in the optimal approximation order).
The unmodified English and Chinese text is supplied in \code{ORIGINAL.md}.

We interpret ``minimal'' as minimizing the number of facets. The phrase
``is'' allows a literal equality, whereas ``optimal approximation order''
suggests an asymptotic statement. We address both, as well as the weaker
positive two-sided order. In each reading a leading or optimal-order
constant is finite, strictly positive, and independent of $\varepsilon$.
It suffices to refute dimension $d=2$, even if a formula is required only
on an arbitrarily small interval $0<\varepsilon<\delta$.

\begin{theorem}\label{thm:main}
There do not exist $a,b,\delta>0$ such that for every
$0<\varepsilon<\delta$ an attained minimum facet count $n$ satisfies
\[
 a\sqrt\varepsilon\leq n\leq b\sqrt\varepsilon.
\]
In particular the printed exact formula and positive leading asymptotic
are false in dimension two, and hence the dimension-uniform conjecture
is false.
\end{theorem}

No sign-corrected rate $\varepsilon^{-(d-1)/2}$ is asserted or refuted
here. A zero constant is not a positive leading or optimal-order
constant. We also make no claim about a bare one-sided upper-bound
formulation. These are distinctions of scope, not alterations of the
printed exponent.

\newpage
\section{Geometric definitions and attainment}
The formal ambient space is \code{EuclideanSpace $\R$ (Fin d)}.
A polytope is a convex hull $P=\conv(S)$ of a finite set of points.
A facet is a nonempty proper exposed face $F$ of $P$ with
\[
 \dim\operatorname{aff}(F)+1=\dim\operatorname{aff}(P).
\]
These are intrinsic dimensions. In the full-dimensional case they give
the usual codimension-one facets. The definition permits all convex
polytopes, without imposing containment in or around the ball, symmetry,
or a prescribed vertex arrangement. Volume is the standard Lebesgue
volume on Euclidean space; symmetric difference is the actual set
operation. For positive $\varepsilon$, the nonnegative extended-real
factor used in Lean is precisely $\varepsilon$.

The relation \code{HasFacetCount P n} means that $\{0,\ldots,n-1\}$ is
in bijection with the set of these geometric facets. It is not a
totalized cardinality of an infinite set. The relation
\code{IsMinimumFacetCount d $\varepsilon$ n} says that an actual
approximating polytope has $n$ facets and every admissible facet count
$m$ satisfies $n\leq m$.

\begin{lemma}\label{lem:finite}
Every polytope has finitely many exposed faces and hence finitely many
facets. Its finite facet count is unique. For each error, an attained
minimum facet count exists if and only if an approximating polytope
exists; that minimum is unique.
\end{lemma}
\begin{proof}
Write $P=\conv(S)$ with $S$ finite. The set $P$ is compact. Every exposed
face $F$ is compact and convex, and its extreme points are among the
extreme points of $P$, which are among $S$. By the finite-dimensional
Krein--Milman theorem,
$F=\overline{\conv(\operatorname{ext}F)}$.
Thus exposed faces belong to the image of the finite power set of $S$
under $A\mapsto\overline{\conv(A)}$. Facets form a subset of this finite
set. Two bijections with finite ordinal sets force their natural-number
cardinalities to agree.

If an approximating polytope exists, its finite facet count makes the
set of admissible natural counts nonempty. Its least element is an
admissible count, so a polytope attaining it exists. Conversely an
attained minimum includes such a polytope by definition. Two minima
satisfy both opposite inequalities, so are equal.
\end{proof}

All assertions in Lemma~\ref{lem:finite} are formalized. The disproof
does not need a separate construction of approximants at every error:
each of the three asserted rates already requires actual minimum
witnesses at sufficiently small errors. We prove that no such witnesses
can satisfy the printed rate, even before using their additional
geometric properties.

\section{The integer obstruction}
Fix arbitrary $a,b,\delta>0$ and set
\[
 t=\tfrac12\min\{\sqrt\delta,1/b\},\qquad \varepsilon=t^2.
\]
Then $t>0$, $t<\sqrt\delta$, and $t<1/b$. Hence
$0<\varepsilon<\delta$, $\sqrt\varepsilon=t$, and $bt<1$.
If $n\in\N$ satisfied the two bounds in Theorem~\ref{thm:main}, then
\[
 0<at\leq n\leq bt<1.
\]
There is no natural number strictly between $0$ and $1$. This proves
Theorem~\ref{thm:main}. The obstruction covers every natural count;
there is no numerical search or finite sampling of errors.

\newpage
\section{The three readings of the conjecture}
For $d=2$ the printed real exponent is $(2-1)/2=1/2$, so the rate is
$\sqrt\varepsilon$ exactly.

\paragraph{Exact formula.}
An equality $n=c\sqrt\varepsilon$ with $c>0$ implies the two-sided bounds
with $a=b=c$, contradicting Theorem~\ref{thm:main}.

\paragraph{Positive leading asymptotic.}
Suppose the minimum counts satisfy
$n(\varepsilon)/(c\sqrt\varepsilon)\longrightarrow1$ as
$\varepsilon\to0^+$, for $c>0$. Choose the tolerance $1/2$ in the
ordinary epsilon-delta definition. At all sufficiently small positive
errors the ratio lies between $1/2$ and $3/2$, and therefore
\[
 (c/2)\sqrt\varepsilon\leq n(\varepsilon)
 \leq (2c)\sqrt\varepsilon.
\]
Theorem~\ref{thm:main} again contradicts the claim. The formal predicate
quantifies actual minimum witnesses relationally for every tolerance;
Lemma~\ref{lem:finite} makes their count unique wherever they exist.
The auxiliary module also proves the corresponding obstruction for
Mathlib's standard filter-based asymptotic equivalence of functions.

\paragraph{Optimal two-sided order.}
An order $\Theta(\sqrt\varepsilon)$ with positive fixed lower and upper
constants is exactly the already-refuted pair of bounds. Thus the
disproof does not rely solely on literal equality of an integer with a
noninteger. It rules out the natural optimal-order interpretation too.
Since no positive constant has the claimed behavior, asking for that
constant to be explicit cannot restore the statement.

\section{Lean correspondence and reproduction}
\code{Geometry.lean} defines the geometric objects and proves finiteness,
unique counts, and the equivalence between minimum attainment and
existence of an approximant. \code{Vanishing.lean} proves the integer
obstruction for every positive choice of constants and error interval,
including the standard asymptotic-equivalence version.

\code{Solution.lean} defines all three readings using attained geometric
minimum witnesses. The final universal negations in namespace
\code{TLMC1451} are:
\begin{flushleft}
\code{conjecture\_exact\_false}\\
\code{conjecture\_asymptotic\_false}\\
\code{conjecture\_order\_false}.
\end{flushleft}
Each follows from a corresponding dimension-two theorem; the universal
domain $d\geq2$ contains that counterexample dimension.

The project pins Lean 4.19.0 and Mathlib v4.19.0, with exact dependency
revisions in its manifest. From the supplied \code{lean/} directory,
use the supplied manifest to obtain dependencies and the Mathlib cache
with \code{lake exe cache get}, then run \code{lake build}.
The supplied verification script additionally replays every authored
module with warnings treated as errors and inspects compiled axiom
dependencies. See \code{VERIFICATION.md} for the recorded commands,
source identities, and audit scope. The mathematical argument is
symbolic; it relies on no auxiliary numerical computation. Local
verification and independent pre-submission review do not constitute
maintainer acceptance.
\end{document}
